\section{Method}
\label{sec:method}

We extract a shared indexator by jointly projecting numerical solutions
of source tasks and comparing the resulting structures after numerical
refitting. Projection proposes groups; prediction error selects the
structure.

\subsection{Source solutions and their alignment}
\label{sec:method-source}

Write $T:=|\mathcal T_{\mathrm{tr}}|$ and
$\boldsymbol\eta_\tau:=(C_{\tau,1},\ldots,C_{\tau,K},
\boldsymbol\theta_\tau)$ for the trainable numerical parameters.
For each source task, fitting uses $S_\tau^{\mathrm s}$.
Its query sample is partitioned into $V_\tau^{\mathrm{num}}$, used
for numerical checkpoints and proposal construction, and
$V_\tau^{\mathrm{str}}$, reserved for the final structure comparison.
Target tasks do not enter this procedure.

First, unconstrained numerical fitting gives
$\widetilde f_\tau=f_{\widetilde{\mathbf w}_\tau,
\widetilde{\boldsymbol\theta}_\tau}$, with
$\widetilde{\mathbf w}_\tau\in\mathbb R^P$.
Dense starts are fitted on support data and selected on
$V_\tau^{\mathrm{num}}$. The numerical updates are specified below;
for dense fitting their variables are the unrestricted network parameters.

\begin{definition}[Source descriptor]
    For a fixed alignment of a source solution, its descriptor is
    \begin{equation}
        \mathbf z_\tau(\mathbf a_\tau)
        =\mathbf B_\tau\mathbf a_\tau+\mathbf v_\tau,
        \qquad
        \mathbf u_\tau(\mathbf a_\tau)\in(0,\infty)^P,
        \label{eq:method-descriptor}
    \end{equation}
    where $\mathbf a_\tau\in\mathbb R^{r_\tau}$ contains continuous
    calibration variables, $\mathbf B_\tau\in\mathbb R^{P\times r_\tau}$
    and $\mathbf v_\tau\in\mathbb R^P$ are obtained from the source
    solution, and $\mathbf u_\tau$ assigns positive discrepancy scales.
\end{definition}

The identity descriptor has $r_\tau=0$,
$\mathbf v_\tau=\widetilde{\mathbf w}_\tau$ and
$\mathbf u_\tau=\mathbf 1$, and is available for any network.
Architecture-specific descriptors can account for parameter symmetries.
Discrete alignments, such as hidden-unit permutations, are fixed for
each projection run; continuous calibration is updated within the run.
Adjacent parameters must be transformed together when an alignment is
claimed to preserve the predictor.

\subsection{Joint discrete projection}
\label{sec:method-projection}

For a binary support $\mathbf M\in\{0,1\}^P$, let
$\mathcal I_K(\mathbf M):=\{\mathbf I\in\mathcal I_K:
I_p=0\Leftrightarrow M_p=0\text{ for all }p\}$.
With $\mathbf c_\tau:=(C_{\tau,1},\ldots,C_{\tau,K})^{\mathsf T}$,
define the lookup matrix by
\begin{equation}
    (\mathbf A_{\mathbf I})_{pj}:=\mathbf 1\{I_p=j\},
    \qquad
    \mathbf w_\tau=\mathbf A_{\mathbf I}\mathbf c_\tau.
    \label{eq:method-lookup-matrix}
\end{equation}
We score a joint projection by
\begin{equation}
    \begin{aligned}
        \mathcal Q(\mathbf I,\{\mathbf c_\tau,\mathbf a_\tau\})
        &:=\frac{1}{TP}\sum_{\tau,p}
        \rho_\nu\!\left(
            \frac{(\mathbf A_{\mathbf I}\mathbf c_\tau)_p
                -z_{\tau,p}(\mathbf a_\tau)}
                 {u_{\tau,p}(\mathbf a_\tau)}\right),\\
        \rho_\nu(v)&:=\nu\log(1+v^2/\nu),\qquad \nu>0.
    \end{aligned}
    \label{eq:joint-projection}
\end{equation}
Assignments are shared across tasks; code values and calibrations are
task-specific. This tests whether positions admit common values
\emph{within} each task.

With $\mathbf I$ and the discrepancy scales fixed, calibration is
a constrained weighted least-squares problem. Set
\begin{equation}
    \begin{gathered}
        \boldsymbol\xi_\tau
            :=(\mathbf c_\tau^{\mathsf T},\mathbf a_\tau^{\mathsf T})^{\mathsf T},
        \qquad
        \mathbf E_{\tau,\mathbf I}:=
            [\mathbf A_{\mathbf I},-\mathbf B_\tau],\\
        \boldsymbol\xi_\tau^+
        \in\operatorname*{arg\,min}_{\mathbf G_{\tau,\mathbf I}
            \boldsymbol\xi=\mathbf h_{\tau,\mathbf I}}
        \left\|\mathbf D_\tau^{-1}
            (\mathbf E_{\tau,\mathbf I}\boldsymbol\xi-\mathbf v_\tau)
        \right\|_2^2,
        \qquad
        \mathbf D_\tau:=\operatorname{diag}
            \mathbf u_\tau(\mathbf a_\tau).
    \end{gathered}
    \label{eq:method-calibration}
\end{equation}
The feasible affine constraints $\mathbf G_{\tau,\mathbf I}\boldsymbol\xi
=\mathbf h_{\tau,\mathbf I}$ fix calibration scale freedoms; they
are absent for the identity descriptor. A minimum-norm minimizer is
used when the solution is nonunique. Equivalently, calibration solves
\begin{equation}
    \begin{pmatrix}
        \mathbf E^{\mathsf T}\mathbf D^{-2}\mathbf E & \mathbf G^{\mathsf T}\\
        \mathbf G & \mathbf 0
    \end{pmatrix}
    \begin{pmatrix}\boldsymbol\xi^+\\\boldsymbol\lambda\end{pmatrix}
    =
    \begin{pmatrix}
        \mathbf E^{\mathsf T}\mathbf D^{-2}\mathbf v\\\mathbf h
    \end{pmatrix},
    \label{eq:method-calibration-system}
\end{equation}
with task and indexator subscripts suppressed. The scales are recomputed
between solves. This least-squares update is a surrogate for the robust
score; monotone decrease of $\mathcal Q$ is not assumed.

\begin{algorithm}[htbp]
    \small
    \DontPrintSemicolon
    \SetAlgoNoEnd
    \KwIn{Aligned descriptors; support $\mathbf M$;
        $\mathbf I^{(0)}\in\mathcal I_K(\mathbf M)$;
        initial calibrations; budgets $N_{\mathrm{proj}},N_{\mathrm{cal}}$.}
    \KwOut{A proposed indexator and its projection score.}
    $\mathbf I\leftarrow\mathbf I^{(0)}$\;
    \For{$k=0,\ldots,N_{\mathrm{proj}}$}{
        \ForEach{$\tau\in\mathcal T_{\mathrm{tr}}$}{
            \For{$q=1,\ldots,N_{\mathrm{cal}}$}{
                $\mathbf D_\tau\leftarrow\operatorname{diag}
                    \mathbf u_\tau(\mathbf a_\tau)$\;
                $\mathbf E_{\tau,\mathbf I}\leftarrow
                    [\mathbf A_{\mathbf I},-\mathbf B_\tau]$\;
                $\boldsymbol\xi_\tau\leftarrow$
                    solution of \eqref{eq:method-calibration-system}\;
            }
        }
        \If{$k=N_{\mathrm{proj}}$}{\textbf{break}\;}
        \ForEach{$p$ with $M_p=1$ and $j=1,\ldots,K$}{
            $d_{pj}\leftarrow\displaystyle\sum_\tau
                \rho_\nu\!\left(
                    \frac{C_{\tau,j}-z_{\tau,p}(\mathbf a_\tau)}
                         {u_{\tau,p}(\mathbf a_\tau)}\right)$\;
        }
        \tcp{The same label must fit every source task.}
        $I_p^+\leftarrow0$ if $M_p=0$; otherwise
            $I_p^+\in\operatorname*{arg\,min}_{j=1,\ldots,K}d_{pj}$
            ($p=1,\ldots,P$)\;
        \If{$\mathbf I^+=\mathbf I$}{\textbf{break}\;}
        $\mathbf I\leftarrow\mathbf I^+$\;
    }
    \KwRet{$(\mathbf I,\mathcal Q(\mathbf I,
        \{\mathbf c_\tau,\mathbf a_\tau\}))$}\;
    \caption{Joint calibration and shared assignments}
    \label{alg:joint-projection}
\end{algorithm}

Assignment ties use a fixed ordering of labels.
Single-source runs of Algorithm~\ref{alg:joint-projection} provide
initial structures. After numerical fitting across tasks, a subset
is retained as anchors $\mathcal A$. For each anchor alignment,
alternative supports include its own mask and source-consensus masks
\begin{equation}
    M_p^{(\lambda,v)}
    :=\mathbf 1\!\left\{
        \sum_\tau\mathbf 1\!\left\{
            |z_{\tau,p}(\mathbf a_\tau^{(0)})|
            >\max\{\lambda u_{\tau,p}(\mathbf a_\tau^{(0)}),\varepsilon\}
        \right\}\ge v\right\},
    \label{eq:method-support-proposals}
\end{equation}
where $\varepsilon>0$ and the finite sets of thresholds $\lambda$
and vote counts $v\in\{1,\ldots,T\}$ are prescribed.
Joint runs start from anchors, single-source assignments and random
active labels.

Let $\mathcal G$ be a finite group of known admissible label and
parameter permutations. We identify duplicates by the canonical form
$\kappa(\mathbf I):=\min_{g\in\mathcal G}^{\mathrm{lex}}g\mathbf I$.
The bank $\mathcal B$ contains the best canonically distinct proposals
ranked by $\mathcal Q$, together with $\kappa(\mathcal A)$.
This removes known symmetries without requiring a test for all
functional equivalences.

\subsection{Numerical refitting and functional selection}
\label{sec:method-selection}

For a fixed candidate, write
$F_{\tau,\mathbf I}(\boldsymbol\eta):=
\mathcal L_{S_\tau^{\mathrm s}}(f_{\mathbf I,\boldsymbol\eta})$.
One numerical fitting rule uses the damped direction and finite line search
\begin{equation}
    \begin{aligned}
        \mathbf g&:=\nabla F_{\tau,\mathbf I}(\boldsymbol\eta),
        \qquad
        \boldsymbol\Delta:=
            -(\mathbf H+\lambda\mathbf{Id})^{-1}\mathbf g,\\
        \mathcal S(\boldsymbol\eta,\boldsymbol\Delta)
        &:=\bigl\{\beta^l:0\le l<N_{\mathrm{ls}},\
            F_{\tau,\mathbf I}(\boldsymbol\eta+\beta^l\boldsymbol\Delta)
            <F_{\tau,\mathbf I}(\boldsymbol\eta)\bigr\},\\
        \alpha&:=\max\bigl(\{0\}\cup
            \mathcal S(\boldsymbol\eta,\boldsymbol\Delta)\bigr),
    \end{aligned}
    \label{eq:method-numerical-update}
\end{equation}
where $\mathbf H=\mathbf H_{\tau,\mathbf I}(\boldsymbol\eta)\succeq0$,
$\lambda>0$, $0<\beta<1$, and $N_{\mathrm{ls}}$ bounds the line search.
For mean squared error, let $\mathbf r$ stack the $n_\tau$ scalar
prediction residuals on support data. The Gauss--Newton choice is
\begin{equation}
    \mathbf J:=\frac{\partial\mathbf r}{\partial\boldsymbol\eta},
    \qquad
    \mathbf H_{\tau,\mathbf I}(\boldsymbol\eta)
        :=\frac{2}{n_\tau}\mathbf J^{\mathsf T}\mathbf J.
    \label{eq:method-gn-matrix}
\end{equation}
The choice $\mathbf H=\mathbf0$ gives a scaled gradient step.
Only independent trainable coordinates are included in $\boldsymbol\eta$.

For each candidate and task, several initial states
$\{\boldsymbol\eta_{\tau,r}^{(0)}\}_{r=1}^{R}$ generate a finite
set of iterates. They can include that task's source-derived
initializations. Numerical states and the shared structure are selected by
\begin{equation}
    \begin{aligned}
        (\widehat r_\tau,\widehat k_\tau)(\mathbf I)
        &\in\operatorname*{arg\,min}_{r,k}
        \mathcal L_{V_\tau^{\mathrm{num}}}
            (f_{\mathbf I,\boldsymbol\eta_{\tau,r}^{(k)}}),\\
        \widehat{\mathbf I}
        &\in\operatorname*{arg\,min}_{\mathbf I\in\mathcal B}
        \frac1T\sum_\tau\mathcal L_{V_\tau^{\mathrm{str}}}
            (f_{\mathbf I,\widehat{\boldsymbol\eta}_\tau(\mathbf I)}).
    \end{aligned}
    \label{eq:method-structure-selection}
\end{equation}
Here $\widehat{\boldsymbol\eta}_\tau(\mathbf I)$ is the iterate
selected by the first line. The bank and these iterates are fixed
before the second line is evaluated.

\begin{algorithm}[htbp]
    \small
    \DontPrintSemicolon
    \SetAlgoNoEnd
    \KwIn{Bank $\mathcal B$; source samples; numerical starts;
        fitting budget $N_{\mathrm{fit}}$.}
    \KwOut{Shared indexator $\widehat{\mathbf I}$.}
    \ForEach{$\mathbf I\in\mathcal B$ and $\tau\in\mathcal T_{\mathrm{tr}}$}{
        $b_\tau(\mathbf I)\leftarrow+\infty$\;
        \For{$r=1,\ldots,R$}{
            $\boldsymbol\eta\leftarrow\boldsymbol\eta_{\tau,r}^{(0)}$\;
            \For{$k=0,\ldots,N_{\mathrm{fit}}$}{
                $v\leftarrow\mathcal L_{V_\tau^{\mathrm{num}}}
                    (f_{\mathbf I,\boldsymbol\eta})$\;
                \If{$v<b_\tau(\mathbf I)$}{
                    $(b_\tau(\mathbf I),\widehat{\boldsymbol\eta}_\tau(\mathbf I))
                        \leftarrow(v,\boldsymbol\eta)$\;
                }
                \If{$k=N_{\mathrm{fit}}$}{\textbf{break}\;}
                $\mathbf g\leftarrow\nabla F_{\tau,\mathbf I}(\boldsymbol\eta)$\;
                $\mathbf H\leftarrow\mathbf H_{\tau,\mathbf I}(\boldsymbol\eta)$\;
                $\boldsymbol\Delta\leftarrow
                    -(\mathbf H+\lambda\mathbf{Id})^{-1}\mathbf g$\;
                $\alpha\leftarrow\max(\{0\}\cup
                    \mathcal S(\boldsymbol\eta,\boldsymbol\Delta))$\;
                \If{$\alpha=0$}{\textbf{break}\;}
                $\boldsymbol\eta\leftarrow
                    \boldsymbol\eta+\alpha\boldsymbol\Delta$\;
            }
        }
    }
    \tcp{Structure-selection data enter only after all numerical choices.}
    $s(\mathbf I)\leftarrow\displaystyle\frac1T\sum_\tau
        \mathcal L_{V_\tau^{\mathrm{str}}}
            (f_{\mathbf I,\widehat{\boldsymbol\eta}_\tau(\mathbf I)})$
        for each $\mathbf I\in\mathcal B$\;
    $\widehat{\mathbf I}\in\operatorname*{arg\,min}_{\mathbf I\in\mathcal B}
        s(\mathbf I)$\;
    \KwRet{$\widehat{\mathbf I}$}\;
    \caption{Numerical adaptation and structure selection}
    \label{alg:structure-selection}
\end{algorithm}

On a target task, $\widehat{\mathbf I}$ is fixed and all its numerical
parameters are initialized afresh and fitted on the target support
sample. Query loss measures transfer. Recovery depends on whether
the finite bank contains a useful structure and whether numerical
fitting can reveal it.

\subsection{ReLU implementation}
\label{sec:method-relu-instance}

Our experimental network is
\begin{equation}
    f_\tau(\mathbf x)
    =H^{-1/2}\mathbf A_\tau^{\mathsf T}
        \rho(\mathbf W_\tau^{\mathsf T}\mathbf x+\mathbf b_\tau)
        +\mathbf d_\tau,
    \label{eq:method-relu-model}
\end{equation}
where $\mathbf W_\tau\in\mathbb R^{D\times H}$,
$\mathbf A_\tau\in\mathbb R^{H\times O}$,
$\mathbf b_\tau\in\mathbb R^H$, $\mathbf d_\tau\in\mathbb R^O$, and
$\rho$ is coordinatewise ReLU. Only $\mathbf W_\tau$ is structured;
the biases and output layer are free.

For synthetic regression, dense solutions are refined by Gauss--Newton
steps and residual-based reinitialization of weak neurons. Normalized columns
$\mathbf r_{\tau,h}:=\widetilde{\mathbf W}_{\tau,:,h}/
\|\widetilde{\mathbf W}_{\tau,:,h}\|_2$ define
\begin{equation}
    z_{\tau,dh}(\mathbf s_\tau)=s_{\tau,h}r_{\tau,dh},
    \qquad
    u_{\tau,dh}(\mathbf s_\tau)
        =\max\{|s_{\tau,h}|,\varepsilon_s\}
            \max\{\widehat u_{\tau,h},\varepsilon_u\},
    \label{eq:method-relu-descriptor}
\end{equation}
with positive floors $\varepsilon_s,\varepsilon_u$.
Reliability weights $\widehat u_{\tau,h}$ are estimated from residuals
and the pseudoinverse of the empirical Jacobian Gram matrix, propagated
through column normalization.

For anchor directions $\overline{\mathbf r}_{\tau,h}$ and support
$\mathbf M^{\mathrm a}$, hidden-unit alignment solves
\begin{equation}
    \begin{aligned}
        \pi_\tau&\in\operatorname*{arg\,min}_{\pi\in\mathfrak S_H}
            \sum_{h=1}^H D_{\tau,h,\pi(h)},\\
        D_{\tau,hk}&:=\sum_d
            \mathbf1\{M^{\mathrm a}_{dh}\ne\widetilde M_{\tau,dk}\}
            +\omega\min_{\sigma\in\{-1,1\}}
                \|\overline{\mathbf r}_{\tau,h}
                    -\sigma\mathbf r_{\tau,k}\|_2^2,
    \end{aligned}
    \label{eq:method-relu-alignment}
\end{equation}
where $\mathfrak S_H$ is the permutation set,
$\widetilde{\mathbf M}_\tau$ is the thresholded source support,
and $\omega>0$ weights direction agreement.
Calibration fixes one column scale per connected component of the
code--neuron incidence graph, connecting each code to the neurons using it.

The implemented signed codebook has $m$ independent values:
\begin{equation}
    \mathbf C_\tau
        =(0,\mathbf c_\tau^{\mathsf T},-\mathbf c_\tau^{\mathsf T})^{\mathsf T},
    \qquad K=2m.
    \label{eq:method-paired-codebook}
\end{equation}
For this restricted numerical family, replace $\mathbf A_{\mathbf I}$
by the signed lookup matrix with row $\sigma\mathbf e_j^{\mathsf T}$
for code $\sigma c_{\tau,j}$, where $\mathbf e_j$ is the $j$th standard
basis vector. Writing its signed labels as
$J_{dh}\in\{0,\pm1,\ldots,\pm m\}$, each reference source $t$ yields
the oriented proposal $J^{(t)}_{dh}=\sigma_{t,h}J_{dh}$, where
$\sigma_{t,h}:=2\mathbf1\{s_{t,h}\ge0\}-1$.
Reversing a ReLU active side is checked by functional fitting.
Canonicalization permutes and changes signs of code pairs jointly.
Empty independent groups are reseeded from groups containing several
positions.

For support inputs $\{\mathbf x_n\}$, define
$\Phi_{nh}:=H^{-1/2}\rho(\mathbf x_n^{\mathsf T}\mathbf W_{\tau,:,h}+b_{\tau,h})$
and let $\mathbf Y$ collect their targets. For squared-loss refinement,
the free output layer is refitted at every iterate and line-search trial by
\begin{equation}
    \begin{pmatrix}\mathbf A_\tau\\\mathbf d_\tau^{\mathsf T}\end{pmatrix}
    \in\operatorname*{arg\,min}_{\mathbf U}
        \left\{\|[\boldsymbol\Phi,\mathbf1]\mathbf U-\mathbf Y\|_{\mathrm F}^2
            +\varepsilon_{\mathrm{lin}}\|\mathbf U\|_{\mathrm F}^2\right\},
    \label{eq:method-free-head}
\end{equation}
with a small stabilizer $\varepsilon_{\mathrm{lin}}>0$.
Hidden biases and codes remain trainable. Gauge normalization uses
$a=\|\mathbf c_\tau\|_2/\sqrt m$ and
$(\mathbf c_\tau,\mathbf b_\tau,\mathbf A_\tau)
\leftarrow(\mathbf c_\tau/a,\mathbf b_\tau/a,a\mathbf A_\tau)$
for $a>0$.
These descriptor, alignment and solver choices specify the ReLU
implementation of the general projection-and-selection procedure.
For the real-data experiments we use Adam numerical fitting and
direction dispersion across independently fitted models as discrepancy
scales. These changes preserve the shared projection and functional
selection stages; their finite search budgets are reported separately.
